\begin{abstract}
For the rail transit service-planning problem, we first analyze the section passenger-flow and OD structure to identify the high-demand overlap zone and generate candidate large-small-route service patterns. On this basis, a combined operating-cost and service-level criterion is used to select the final operation plan, and a timetable is then generated through the corresponding arrival-departure recurrence constraints. The obtained plan therefore does not stop at a route-pattern comparison, but yields a machine-readable timetable together with the selected service strategy. Furthermore, capacity, tracking, and dwell-time audits are introduced to examine the operational feasibility of the generated timetable, while scenario analysis is used to summarize the cost-service tradeoff under demand or interval perturbation. The resulting route closes the full chain of flow reading, service-pattern selection, timetable generation, feasibility audit, and operating recommendation, and thus provides a reusable research template for timetable-type mathematical modeling papers.

\keywords{mathematical modeling\quad rail transit\quad service planning\quad timetable generation\quad feasibility audit}
\end{abstract}
